%%=============================================================================
%% CAPITULO 2 -- REFERENCIAL TEORICO
%% Versao dummy: resumo do capitulo 02-referencial-teorico da dissertacao real.
%%=============================================================================

\section{Referencial teórico}

A técnica de RAG enriquece o contexto do modelo de linguagem com documentos
recuperados de uma base de conhecimento indexada, o que reduz o risco de
alucinações e atualiza o conhecimento do modelo sem novo treinamento, conforme
proposto por \citeonline{lewis2020}. Nesta pesquisa, a RAG se combina com o
acoplamento, via \textit{function calling}, de um módulo de cálculo
quantitativo determinístico. A Figura~\ref{fig:arquitetura} sintetiza a
arquitetura proposta e mostra a segregação entre o processamento numérico
determinístico e a geração de linguagem natural.

\begin{figure}[htb]
  \caption{Arquitetura de agente único com segregação entre cálculo determinístico e geração de linguagem}
  \label{fig:arquitetura}
  \centering
  \footnotesize
  \begin{tikzpicture}[
    node distance=6mm,
    box/.style={draw, rounded corners=2pt, align=center, inner sep=4pt, minimum height=8mm},
    ext/.style={box, dashed, minimum width=30mm},
    core/.style={box, thick, minimum width=34mm},
    io/.style={box, minimum width=30mm},
    arrow/.style={-{Latex[length=2mm]}, thick},
    feed/.style={-{Latex[length=2mm]}, dashed}
  ]
    \node[io] (cliente) {Cliente pessoa física\\{\scriptsize (pergunta em linguagem natural)}};
    \node[core, below=of cliente] (agente) {\textbf{Agente LLM único}\\{\scriptsize orquestração via \textit{function calling}}};

    \node[core, below left=10mm and 2mm of agente] (python) {Módulo determinístico\\Python\\{\scriptsize RFL, IR/IOF, \textit{duration}}};
    \node[core, below right=10mm and 2mm of agente] (rag) {Base RAG normativa\\{\scriptsize CVM, CMN, Tesouro, IR/IOF}};

    \node[ext, below=9mm of python] (apis) {Dados públicos\\{\scriptsize BCB e Tesouro Transparente}};
    \node[ext, below=9mm of rag] (normas) {Normativos vigentes\\{\scriptsize com metadados de vigência}};

    \node[io, above right=6mm and 4mm of cliente] (xai) {Log estruturado XAI\\{\scriptsize trilha de auditoria}};

    \draw[arrow] (cliente) -- node[midway, left, xshift=-1mm] {\scriptsize pergunta} (agente);
    \draw[arrow] (agente) -- (python) node[midway, left, xshift=-1mm] {\scriptsize cálculo};
    \draw[arrow] (agente) -- (rag) node[midway, right, xshift=1mm] {\scriptsize consulta};
    \draw[arrow] (python) -- (agente);
    \draw[arrow] (rag) -- (agente);
    \draw[feed] (apis) -- (python);
    \draw[feed] (normas) -- (rag);
    \draw[arrow] (agente.west) to[out=180, in=180] node[midway, left, align=right] {\scriptsize resposta em\\\scriptsize linguagem natural} (cliente.west);
    \draw[arrow] (agente.east) to[out=0, in=-90] node[midway, right, yshift=-2mm] {\scriptsize log} (xai.south);
  \end{tikzpicture}
  \fonte{Elaborada pelo autor (2026).}
\end{figure}

No domínio financeiro-regulatório brasileiro, a aplicação de RAG sobre os
normativos da Comissão de Valores Mobiliários (CVM), as resoluções do Conselho
Monetário Nacional (CMN) e a legislação tributária permite ao agente responder
com fundamentação documental rastreável. Esse é um requisito essencial para a
conformidade com a Resolução CVM nº~30/2021, que estabelece o dever de
verificação da adequação dos produtos, serviços e operações ao perfil do
cliente (\textit{suitability}), nos seguintes termos:

\begin{citacao}
Art. 2º As pessoas habilitadas a atuar como integrantes do sistema de
distribuição e os consultores de valores mobiliários não podem recomendar
produtos, realizar operações ou prestar serviços sem que verifiquem sua
adequação ao perfil do cliente. [...] Art. 3º As pessoas referidas no art. 2º
devem verificar se: I – o produto, serviço ou operação é adequado aos
objetivos de investimento do cliente; II – a situação financeira do cliente é
compatível com o produto, serviço ou operação; e III – o cliente possui
conhecimento necessário para compreender os riscos relacionados ao produto,
serviço ou operação. \parencite[p.~1]{resolucaocvm30}
\end{citacao}
